% Same Terminal Signal, Different Action
% Evidence-bounded preprint generated from frozen repository artifacts.
\pdfsuppressptexinfo=-1
\pdfinfoomitdate=1
\pdftrailerid{}

\documentclass{article}

\usepackage[preprint,nonatbib]{neurips_2026}
\usepackage[numbers,sort&compress]{natbib}
\usepackage[utf8]{inputenc}
\usepackage[T1]{fontenc}
\usepackage{amsmath,amssymb,amsfonts,amsthm}
\usepackage{booktabs}
\usepackage{enumitem}
\usepackage{microtype}
\usepackage{multirow}
\usepackage{tabularx}
\usepackage{xcolor}
\usepackage{xspace}
\usepackage{hyperref}
\usepackage{url}

\hypersetup{
  colorlinks=true,
  linkcolor=blue!55!black,
  citecolor=blue!55!black,
  urlcolor=blue!55!black
}

\newtheorem{proposition}{Proposition}
\newtheorem{corollary}{Corollary}
\newtheorem{remark}{Remark}
\newcommand{\grpo}{\textsc{GRPO}\xspace}
\newcommand{\zvf}{\textsc{ZVF}\xspace}
\newcommand{\E}{\mathbb{E}}
\newcommand{\Prob}{\mathbb{P}}
\newcommand{\ind}{\mathbf{1}}
\newcommand{\method}{\textsc{FailClosed-Conformance}\xspace}

\title{Same Terminal Signal, Different Action:\\
Fail-Closed Conformance and Decision Insufficiency in Group-Relative RL}

\author{Arvind C R\\
PES University\\
\texttt{arvindcr4@gmail.com}}

\begin{document}
\maketitle

\begin{abstract}
Group-relative reinforcement learning can assign exactly zero centered
advantage to a group whose rewards are all equal.  The event is easy to detect
but ambiguous: an all-zero group may reflect an impossible prompt, a capable
policy that was unlucky, a false-negative verifier, or an execution failure.
Those states can require opposite actions even though they expose the same
terminal reward history.  At the same time, nominally identical objectives can
differ across implementations in standardization, row selection, and token
reduction, making an apparent algorithmic result inseparable from software
semantics.

We present an evidence-bounded study of these two problems.  First, we define a
fail-closed conformance chain from objective equation to gradient relation,
adaptive decision, and training record.  A frozen CPU differential contains 14
intended cases for each of TRL and verl and a shared 36-case controller matrix,
while retaining native differences rather than relabeling them as conformance.
Second, we give cost-accounting corollaries of the known binary-starvation
model.  For fixed $G$, increasing group size reduces the expected number of
groups needed to see a contrast; for every $G$, complete-group sampling costs
at least $1/\epsilon$ expected rollouts at success probability $\epsilon$.  We
also give a restricted two-state witness: when $\Delta>0$, the two actions
\{stop, commit to $m$ batch rollouts\} reverse rank across compatible latent
states and incur positive two-point minimax regret.

We then report the outcome of a frozen A100 screening campaign without filling
in missing evidence.  Six scientific units completed in the balanced
equal-length negative control.  Each contained 100 replay steps, 62--65 of
which store exact joint-zero gradient diagnostics; all remaining intended-versus-native
cosines were at least 0.999844.  High cosine was not the registered gate: the
sole completed intended-full cell satisfies the joint cosine/relative-$\ell_2$
criterion on 69/100 steps, below the required 95/100.  The filtered
variable-length positive control was contract-infeasible, four seed-37 units
remain pending for accelerator quota, and no confirmatory matrix ran.  Consequently, this registered
feasibility outcome does not establish a causal training effect or a scientific
negative result about objective semantics.  The paper's contribution is an
auditable boundary between conformance evidence, mathematical necessity, and a
registered feasibility outcome.
\end{abstract}

\section{Introduction}

Reinforcement learning from verifiable rewards (RLVR) commonly samples several
responses to a prompt and centers each reward against its group.  This design
avoids a learned value model and underlies Group Relative Policy Optimization
(GRPO) \citep{shao2024deepseekmath}.  It also has an exact failure mode.  When
all group rewards agree, every centered advantage is zero and the corresponding
score-function update vanishes.  Recent work names and measures this phenomenon
as gradient starvation or advantage collapse and develops direct mitigations
\citep{nie2026gradient,he2026advantage}.  We therefore treat the zero-variance
event as background, not as a new discovery.

The open problem is what a system should infer from the event.  The same vector
of zero rewards can be generated by a prompt with zero success probability, a
prompt with small but nonzero success probability, a verifier that rejects
valid answers, or a broken runtime.  A controller that sees only terminal
outcomes cannot in general distinguish these causes.  Retrying is wasteful in
the first case but valuable in the second; additional verification is
appropriate in the third; repairing the execution path dominates either in the
fourth.  TinyV directly studies verifier false negatives in RL
\citep{xu2025tinyv}; VerifyBench documents broader limitations of
reference-based reward systems \citep{yan2025verifybench}.

A second ambiguity arises below the mathematical specification.  Open RL stacks
expose objectives with similar names but can differ in reward standardization,
which samples are retained, whether reduction is per completion or over all
active tokens, and how importance ratios are clipped.  DAPO, for example,
combines dynamic sampling with a particular token-level loss design
\citep{yu2025dapo}; verl is built on the HybridFlow execution framework
\citep{sheng2024hybridflow}.  A run that terminates successfully need not
implement the intervention that its label suggests.

This paper asks one central question:

\begin{quote}
When a group-relative learner exposes the same terminal starvation statistic,
what evidence is sufficient to choose an intervention, and how can we verify
that the software realizes that intervention before interpreting training
outcomes?
\end{quote}

We make four contributions.

\begin{enumerate}[leftmargin=*,itemsep=2pt]
  \item We define a fail-closed conformance protocol that preserves native
  semantic differences, represents zero-gradient relations without fabricated
  angles, and binds accepted results to source and artifact receipts.
  \item We derive scoped cost-accounting corollaries: exact first-contrast
  reliability, a fixed-$G$ boundary expansion, and an all-$G$ rollout lower
  bound.  These are elementary consequences of the known starvation model,
  not a competing starvation theory.
  \item We prove a two-action batch-retry witness and state the standard
  conditional value-of-information inequality under explicit assumptions.
  \item We release an audited registered-feasibility postmortem.  We explicitly
  do not claim a training effect or controller improvement.
\end{enumerate}

\paragraph{Claim discipline.}
Every substantive statement is assigned one of four statuses: proved,
artifact-verified, externally sourced, or proposed.  The portable companion
\path{verify_claims.py} runs executable mathematical sanity checks and verifies
receipt hashes and internal invariants; it does not replace the analytic proofs
or recompute gradients from checkpoints.  The separate \path{CLAIM_AUDIT.md}
maps claims to exact evidence in a content-addressed review bundle.

\section{Background and Related Work}
\label{sec:background}

\subsection{Group-relative objectives and homogeneous rewards}

For prompt $x$, let $G\geq2$ responses receive rewards
$R_1,\ldots,R_G$.  A representative group-relative advantage is
\begin{equation}
  A_i = \frac{R_i-\bar R}{s_R+\varepsilon},
  \qquad
  \bar R=\frac{1}{G}\sum_{j=1}^{G}R_j .
  \label{eq:group-advantage}
\end{equation}
If all rewards are equal, every numerator is zero.  We use \emph{zero-variance
fraction} (ZVF) only as a descriptive name for the observed fraction of such
groups.  It is not an estimate of a prompt's latent success probability.

For independent binary rewards with success probability $p$, the exact
homogeneous-group probability is
\begin{equation}
  f_G(p)=p^G+(1-p)^G.
  \label{eq:degeneracy}
\end{equation}
Both Eq.~\eqref{eq:degeneracy} and the Jensen amplification caused by prompt
heterogeneity are established prior results \citep{nie2026gradient}.  AVSPO
measures advantage collapse and injects virtual samples to learn from
homogeneous groups \citep{he2026advantage}.  Our novelty claim is not that these
events exist or that group size should adapt.

\subsection{Conformance before causal interpretation}

We use \emph{semantic conformance} to mean agreement with an explicit reference
at four levels:
\begin{equation}
  \text{equation}\longrightarrow\text{loss/gradient}
  \longrightarrow\text{decision}\longrightarrow\text{training record}.
  \label{eq:chain}
\end{equation}
An implementation may intentionally define different semantics.  A conformance
test should report that difference, not call either implementation incorrect by
default.  Causal attribution requires one further step: the material semantic
difference must survive matched data, model state, and compute accounting and
must change a preregistered learning outcome.

This separation follows reproducibility practice in machine learning
\citep{pineau2021reproducibility} but targets the objective-to-gradient boundary
that is especially fragile in LLM RL.  It also makes a negative result useful:
if stacks agree in the target regime, the suite becomes infrastructure rather
than evidence for a new algorithm.

The protocol is also related to differential testing, where disagreements
between independently implemented systems expose semantic faults
\citep{pei2017deepxplore}.  Our narrower focus is not generic model-output
coverage: the compared object is a frozen scalar objective and its gradient,
and every verdict is bound to source and fixture hashes.  This complements
experiment-lineage systems such as MLflow \citep{zaharia2018mlflow}; lineage
records what ran, whereas conformance also asks whether the recorded operation
matches a declared equation.

\section{A Fail-Closed Conformance Protocol}
\label{sec:method}

\subsection{Reference semantics and frozen cases}

The repository's S1 harness evaluates frozen reward and mask fixtures without
loading a model or running an optimizer.  A framework-neutral reference emits
advantages, selected rows, losses, and decision outputs.  Separate adapters
evaluate pinned TRL and verl distributions and record package versions, source
paths, and hashes.  The combined receipt has SHA-256 fixture digest
\texttt{c35916cf...af8ae9b}; it contains 14 intended cases per stack and 36
controller cases with no recorded errors.

Here is the complete scalar objective checked by S1.  Let $\mathcal S$ be the
selected completion rows and $M_{it}\in\{0,1\}$ the completion mask.  Within
each reward group, GRPO, DAPO, and GSPO use sample standard deviation
(Bessel correction); if it is at most $10^{-8}$, $A_i=0$.  DrGRPO uses only the
centered reward.  AERO replaces the zero-variance branch with its frozen
Beta-smoothed posterior scale.  For token log probabilities $\ell_{it}$ and
old log probabilities $\ell^{\rm old}_{it}$, the exact advantage rule is
\begin{equation}
A_i=\begin{cases}
(R_i-\bar R_g)/s_g, & \text{arm}\ne\text{DrGRPO},\ s_g>10^{-8},\\
R_i-\bar R_g, & \text{arm}=\text{DrGRPO},\\
(R_i-\tilde p_g)/\sqrt{\tilde p_g(1-\tilde p_g)},
  & \text{arm}=\text{AERO},\ s_g\leq10^{-8},\\
0, & \text{otherwise},
\end{cases}
\quad \tilde p_g=\frac{k_g+1}{n_g+2},
\label{eq:s1-advantage}
\end{equation}
where $\bar R_g$ and $s_g$ use the active rows of group $g$, and $(k_g,n_g)$
are the frozen AERO success/observation counts.  The loss is
\begin{align}
 \rho_{it}&=\exp(\ell_{it}-\ell^{\rm old}_{it}), \nonumber\\
 h_{it}&=-\min\!\left\{\rho_{it}A_i,
     \operatorname{clip}(\rho_{it};0.8,1+u)A_i\right\}, \nonumber\\
 L_{\rm ref}&=\frac{1}{|\mathcal S|}\sum_{i\in\mathcal S}
   \frac{\sum_t M_{it}h_{it}}{\sum_t M_{it}},
 \label{eq:s1-objective}
\end{align}
where $u=0.28$ for DAPO and $0.20$ otherwise.  GSPO replaces the token
ratio by the sequence ratio
$\exp(\sum_tM_{it}(\ell_{it}-\ell^{\rm old}_{it})/\sum_tM_{it})$.
There is no KL or other auxiliary term in this CPU differential.

The four campaign conditions use the DAPO loss structure in
Eq.~\eqref{eq:s1-objective}, but its training implementation has a separate
$10^{-12}$ zero-standard-deviation floor (S1 uses $10^{-8}$ for its float64
conformance comparisons).  \texttt{intended\_full} standardizes selected
rewards without an additive epsilon, emits zero advantages at that campaign
floor, and averages completion means.  \texttt{epsilon\_only} instead adds
$10^{-4}$ to a nonzero selected standard deviation.
\texttt{reduction\_only} keeps the intended advantages but uses one global
masked-token mean.  \texttt{native\_trl} standardizes all eight reward rows
with $s+10^{-4}$ before selection and uses the global masked-token mean; it does
not use the campaign's explicit floor branch.  Inactive rows have all-zero
masks.  The verifier reads and asserts both source constants.

The harness intentionally separates two questions:
\begin{enumerate}[leftmargin=*,itemsep=2pt]
  \item Does the intended integration reproduce the canonical reference within
  absolute tolerance $10^{-8}$ and relative tolerance $10^{-6}$ in float64?
  \item What does the stack's native objective do on the same fixture?
\end{enumerate}
The combined receipt is \texttt{S1\_PASS} for the first question.  For the
second, it preserves the verdict vectors in Table~\ref{tab:s1}; these verdicts
mean ``different from this reference,'' not ``framework bug.''

\begin{table}[t]
  \centering
  \caption{Frozen S1 conformance receipt.  ``Material'' and ``not tested'' are
  native-versus-reference verdicts; intended integrations pass separately.}
  \label{tab:s1}
  \begin{tabular}{lrrr}
    \toprule
    Stack & Intended cases & Native material & Native not tested \\
    \midrule
    TRL 1.2.0 & 14 & 4 & 1 \\
    verl 0.3.0.post1 & 14 & 1 & 4 \\
    \bottomrule
  \end{tabular}
\end{table}

The controller matrix is shared across stacks: six policies crossed with six
observations.  The policies are fixed $G\in\{8,16\}$, symmetric-ZVF,
failure-only, boundary-aware, and full-triage; the observations are all-wrong,
all-correct, mixed, noisy, missing, and delayed.  Every policy returns
\texttt{recheck} for an unresolved reward.  Otherwise, static policies keep
their named size; symmetric-ZVF escalates any zero-variance group to 16;
failure-only escalates all-wrong groups; boundary-aware escalates all-wrong and
drops all-correct groups; and full-triage escalates all-wrong when the Wilson
upper bound is at least $0.05$ (dropping it otherwise), drops all-correct when
the Wilson lower bound is at least $0.95$, and keeps $G=8$ otherwise.
Appendix~\ref{app:cases} lists every objective case and native verdict; the
review bundle contains exact inputs and outputs.

\subsection{Zero gradients are relations, not angles}

Let $g_a$ and $g_b$ be gradients under two semantics.  A cosine is meaningful
only when both norms are positive.  Here each gradient is the float64 CPU
gradient of the complete selected policy objective in
Eq.~\eqref{eq:s1-objective}; ``zero'' means a bit-exact stored norm of zero,
not a tunable threshold.  Receipts retain both norms and use four logical cases:
\begin{equation}
\operatorname{rel}(g_a,g_b)=
\begin{cases}
  \texttt{joint\_zero}, & \lVert g_a\rVert=\lVert g_b\rVert=0,\\
  \texttt{a\_zero}, & \lVert g_a\rVert=0<\lVert g_b\rVert,\\
  \texttt{b\_zero}, & \lVert g_b\rVert=0<\lVert g_a\rVert,\\
  \texttt{nonzero}, & \lVert g_a\rVert\lVert g_b\rVert>0.
\end{cases}
\label{eq:gradient-relation}
\end{equation}
Cosine and relative $\ell_2$ distance are null in the first three cases.
Joint zero counts as equivalence; a one-sided zero is a categorical material
divergence.  This classification alone makes no general claim about parameter
updates under momentum, weight decay, auxiliary losses, or accumulation.  In
the frozen campaign specifically, $\beta=0$, weight decay is zero, gradient
accumulation is one, and a selected zero gradient explicitly skips
\texttt{optimizer.step}; parameters and AdamW state are unchanged while the
linear scheduler advances.  This avoids both a division-by-zero artifact and
selective deletion of hard replay groups.

\subsection{Failure boundaries and immutable evidence}

The protocol rejects a unit when a required source hash, corpus fingerprint,
package pin, profiler phase, gradient relation, checkpoint manifest, or matched
budget field is absent or inconsistent.  Infrastructure attempts remain
separate from scientific units.  Corpus generation, diagnostic backward
passes, and held-out evaluation are charged explicitly.  A retry may resume
only the greatest checkpoint whose manifest and content fingerprints match.

Development logs record several infrastructure and diagnostic corrections, but
the paper does not treat those anecdotes as scientific evidence.  The claim
boundary begins at immutable accepted receipts: failed attempts remain separate
from accepted units, and the positive-control construction failure remains a
registered feasibility outcome rather than being pooled with later revisions.

\section{What Homogeneous Outcomes Do and Do Not Identify}
\label{sec:theory}

\subsection{Prompt heterogeneity is background, not an estimator}

For $G\geq2$, Eq.~\eqref{eq:degeneracy} is strictly convex on $(0,1)$ because
\begin{equation}
  f_G''(p)=G(G-1)\left[p^{G-2}+(1-p)^{G-2}\right]>0.
\end{equation}
Thus $\E[f_G(P)]\geq f_G(\E[P])$ for a heterogeneous prompt population, with
strict inequality when $P$ is non-degenerate, under the conditional-i.i.d.
binary model.  This does not permit inversion of one observed homogeneous group
into a latent $p$.

\subsection{Time to the first informative group}

Fix a policy and verifier.  Within every group, sample $G$ i.i.d. binary
outcomes with success probability $p$; assume groups are independent and that
$p$ remains stationary until the first contrast.  Define
\begin{equation}
  q_G(p)=1-f_G(p)=1-p^G-(1-p)^G,
\end{equation}
the probability that a group contains both outcomes.

\begin{proposition}[First-contrast reliability]
If $0<p<1$ and $0<\delta<1$, the number $T$ of groups up to and including the
first mixed group is geometric with parameter $q_G(p)$.  Therefore
\begin{equation}
  \E[T]=\frac{1}{q_G(p)},\qquad
  \Prob(T>n)=f_G(p)^n,
\end{equation}
and the smallest number of groups giving probability at least $1-\delta$ of a
contrast is
\begin{equation}
  n_\delta=\left\lceil\frac{\log\delta}{\log f_G(p)}\right\rceil.
  \label{eq:first-contrast}
\end{equation}
At $p\in\{0,1\}$ no finite $n$ suffices.
\end{proposition}

\begin{proof}
Each independent group is mixed with probability $q_G(p)$.  The geometric
mass and tail follow immediately.  Solving $f_G(p)^n\leq\delta$ gives
Eq.~\eqref{eq:first-contrast}; because $0<f_G(p)<1$, division by
$\log f_G(p)<0$ reverses the inequality.
\end{proof}

\begin{proposition}[Fixed-$G$ boundary cost and all-$G$ lower bound]
Fix $G\geq2$ and let $p=\epsilon$ with $\epsilon\downarrow0$ (the case
$p=1-\epsilon$ is symmetric).  Then
\begin{equation}
 q_G(\epsilon)=G\epsilon+O_G(\epsilon^2),\qquad
 \E[\text{groups to contrast}]=\frac{1}{G\epsilon}+O_G(1),
\end{equation}
while the expected number of sampled rollouts is
\begin{equation}
 G\E[T]=\frac{1}{\epsilon}+O_G(1).
 \label{eq:boundary-cost}
\end{equation}
Moreover, for every $G\geq2$ and $0<\epsilon<1$,
\begin{equation}
 q_G(\epsilon)\leq1-(1-\epsilon)^G\leq G\epsilon,
 \qquad G\E[T]=\frac{G}{q_G(\epsilon)}\geq\frac{1}{\epsilon}.
 \label{eq:all-g-cost}
\end{equation}
Thus complete fixed-size groups cannot beat the rare-success waiting-cost lower
bound; the asymptotic coefficient in Eq.~\eqref{eq:boundary-cost} is asserted
only for fixed $G$.
\end{proposition}

\begin{proof}
Use $(1-\epsilon)^G=1-G\epsilon+\binom{G}{2}\epsilon^2+O(\epsilon^3)$ and
$\epsilon^G=O(\epsilon^2)$ for fixed $G\geq2$, then invert the resulting
series.  Multiplication by $G$ yields Eq.~\eqref{eq:boundary-cost}.  For
Eq.~\eqref{eq:all-g-cost}, a mixed group is a subset of the event ``at least one
success,'' and the union bound gives the second inequality.
\end{proof}

This proposition is deliberately narrower than a claim that group size never
helps.  Larger groups can change batching, gradient variance, and optimization
dynamics.  It says only that the first-order sampling burden of discovering a
rare binary outcome remains $1/\epsilon$ for fixed $G$.  The exact lower bound
holds for all $G$, but does not analyze adaptive group sizes, early stopping, or
reuse of partial groups.

\subsection{The same history can require opposite actions}

Consider a controller that has observed $n$ failures from one root prompt.  It
chooses exactly one of two actions: \emph{stop}, or commit in advance to a
batch of $m\geq1$ additional rollouts.  All $m$ costs are paid even if an early
rollout succeeds.  Each rollout costs $c>0$, and observing at least one success
yields benefit $B>0$.
The controller faces one of two latent environments:
\begin{align*}
  E_0 &: p=0,\\
  E_q &: p=q,\quad 0<q<1.
\end{align*}
Both environments assign positive probability to the observed all-failure
history for every finite $n$.

\begin{proposition}[Two-action batch-retry insufficiency]
Suppose
\begin{equation}
  \Delta=B\left[1-(1-q)^m\right]-cm>0.
  \label{eq:delta}
\end{equation}
Then stop is uniquely optimal in $E_0$, while the batch is uniquely optimal in
$E_q$.  Any randomized rule over these two actions that chooses the batch with
probability $\rho$ has two-point worst-case regret
\begin{equation}
  \min_{0\leq\rho\leq1}\max\{\rho cm,(1-\rho)\Delta\}
  =\frac{cm\Delta}{cm+\Delta}>0.
  \label{eq:minimax}
\end{equation}
\end{proposition}

\begin{proof}
Relative to stopping, retry has value $-cm$ in $E_0$ and $\Delta$ in $E_q$.
The same observed history therefore induces opposite optimal actions.  A policy
retrying with probability $\rho$ incurs regret $\rho cm$ in $E_0$ and
$(1-\rho)\Delta$ in $E_q$.  Equalizing these two linear functions gives
$\rho=\Delta/(cm+\Delta)$ and Eq.~\eqref{eq:minimax}.
\end{proof}

The prior $n$ failures establish only that both latent environments are
compatible; they do not enter a posterior calculation.  The result is a
restricted two-action witness, not a theorem about sequential stopping, choosing
$m$, changing verification, or every possible outcome-only controller.  It
shows that terminal outcomes alone need not identify the better action.

\subsection{When side information has value}

Let $O$ denote the terminal-outcome history, $Z$ an additional trace, and
$U(a,\theta)$ utility under action $a$ and latent state $\theta$.  For a finite
action set, the conditional gross Bayes value of observing $Z$ is nonnegative:
\begin{equation}
  \E\!\left[\max_a \E[U(a,\theta)\mid O,Z]\mid O\right]
  \geq
  \max_a \E[U(a,\theta)\mid O].
  \label{eq:value-information}
\end{equation}
This follows because a policy with $Z$ may ignore it and by the tower property.
Equality holds if one action is posterior-optimal for almost every $Z$
conditional on $O$; strict gross value requires a positive expected gap and no
such common action.  Net value subtracts trace-acquisition and downstream action
costs.  Thus ``more telemetry'' is not the result; action-relevant information
under matched total cost is.

\section{Registered Feasibility Outcome and Protocol Postmortem}
\label{sec:campaign}

\subsection{Preregistered design}

The frozen campaign asked whether native-versus-intended objective semantics
change normalized gradient direction and 100-step learning under matched
generated-token and FLOP budgets.  It pinned TRL 1.2.0, Transformers 5.5.4,
Qwen3-1.7B at a specific revision, A100 execution, group size eight, and four
conditions: \texttt{intended\_full}, \texttt{native\_trl},
\texttt{epsilon\_only}, and \texttt{reduction\_only}.  Three screening seeds
were planned in two regimes:

\begin{itemize}[leftmargin=*,itemsep=2pt]
  \item \textbf{balanced equal length}: GSM8K binary rewards
  \citep{cobbe2021gsm8k}, all eight rows active; negative control;
  \item \textbf{filtered variable length}: a MATH-500 slice derived from MATH
  \citep{hendrycks2021math}, select six of sixteen candidates with active
  completion-length coefficient of variation (CV) at least 0.35 and retain two
  inactive rows; positive control.
\end{itemize}

The main hypothesis required equivalence in the balanced regime and material
gradient/learning divergence in the filtered regime.  Its falsifiers included
equivalent gradients in the positive control, inconsistent learning-effect
direction across all three seeds, failure of the causal ablations, or
unverifiable matched compute.

\subsection{Disposition of all planned work}

Table~\ref{tab:state} is recomputed from the current supervisor state.  Ten
accepted jobs comprise three balanced corpora, one non-scientific A100 stack
smoke, and six scientific units.  Acceptance here is a provenance state, not a
claim that the scientific gate passed.

\begin{table}[t]
  \centering
  \caption{Revision r4-2 job disposition on 2026-07-27.}
  \label{tab:state}
  \begin{tabular}{lrp{0.54\linewidth}}
    \toprule
    State & Count & Interpretation \\
    \midrule
    Accepted & 10 & 3 corpora, 1 smoke, 6 scientific units \\
    Contract-infeasible & 14 & Filtered descendants cannot satisfy frozen CV gate \\
    Failed validation & 1 & Filtered seed-11 corpus validation failure \\
    Failed infrastructure & 2 & Preserved, excluded attempts \\
    Pending quota reset & 4 & Balanced seed-37 scientific conditions \\
    \bottomrule
  \end{tabular}
\end{table}

The filtered seed-11 generator failed the frozen contract with observed
completion-length CV $0$, below the required $0.35$, under the 512-token cap.
The intended positive-control variation therefore could not be constructed
without changing the registered treatment.  We descoped that branch rather
than weakening the gate.  The four balanced seed-37 scientific units remain
pending because a fresh A100 allocation probe was rejected for
quota/entitlement.  Neither partial progress nor pending work is counted as a
scientific outcome.

\subsection{Completed scientific units}

Table~\ref{tab:runs} reports all and only final scientific acceptance records.
For a given seed, conditions consume the same immutable ordered corpus and
initial adapter state, as checked by corpus fingerprints and the shared step-0
evaluation hash.  Seed 23 has all four conditions; seed 11 has two.  Causal or
between-condition inference is invalid for this incomplete, unbalanced
screening subset; the displayed quantities are descriptive receipt summaries.

Evaluation used GSM8K test rows 0--127 ($N=128$), greedy decoding
(temperature zero), and the frozen strict final \texttt{\#\#\#\#} integer
parser.  The same rows were evaluated at steps 0, 20, 40, 60, 80, and 100.
Every unit began at $20/128$.  Receipts contain the intermediate curve and the
full token/FLOP ledger; Table~\ref{tab:runs} exposes the endpoint denominator
and charged generation budget.

\begin{table}[t]
  \centering
  \caption{Accepted balanced equal-length units.  Accuracy is count/128;
  every start is 20/128.  Each unit has 100 gradient receipts.  Cosine and
  relative-$\ell_2$ extrema use only nonzero intended/native pairs.}
  \label{tab:runs}
  \scriptsize
  \begin{tabular}{lrrrrrrr}
    \toprule
    Condition & Seed & Start & Final & Joint 0 & Min cos & Max rel-$\ell_2$ & Tokens \\
    \midrule
    epsilon only & 11 & 20 & 17 & 62 & 0.999858 & 0.01684 & 396672 \\
    reduction only & 11 & 20 & 19 & 62 & 0.999844 & 0.01768 & 396672 \\
    epsilon only & 23 & 20 & 17 & 65 & 0.999894 & 0.01458 & 400448 \\
    intended full & 23 & 20 & 20 & 65 & 0.999881 & 0.01544 & 400448 \\
    native TRL & 23 & 20 & 17 & 65 & 0.999894 & 0.01458 & 400448 \\
    reduction only & 23 & 20 & 20 & 65 & 0.999881 & 0.01544 & 400448 \\
    \bottomrule
  \end{tabular}
\end{table}

There were no one-sided-zero intended-versus-native relations in these six
units.  The acceptance records contain 62 joint-zero steps for seed 11 and 65
for seed 23; every such receipt stores both norms as exact zero and leaves
cosine and relative $\ell_2$ null.  The verifier checks those stored invariants
but does not regenerate gradients from private model checkpoints.

For the sole completed intended-full balanced cell (seed 23), 69/100 receipts
satisfy the frozen joint criterion, below the preregistered 95/100 requirement.
Sixty-five are joint-zero; only four of the 35 nonzero receipts satisfy both
cosine $\geq 0.999$ and relative $\ell_2\leq 0.01$.  The very high cosine
extrema are therefore descriptive near-collinearity, not registered
equivalence.  The incomplete screening matrix prevents a full campaign
verdict, but no missing cell can reverse this accepted cell's failure.  The
central causal hypothesis also required the missing filtered positive control
and direction-consistent effects across all three seeds.  Differences in the
six final accuracies are neither causal effects nor evidence of superiority.
No confirmatory seeds ran.

\subsection{What the feasibility outcome establishes}

The evidence supports the following limited claims:
\begin{enumerate}[leftmargin=*,itemsep=2pt]
  \item fail-closed receipts can represent recorded exact-zero gradient pairs
  without dropping data or fabricating cosine similarity;
  \item the balanced records show descriptive near-collinearity, while the sole
  completed intended-full cell fails the registered 95/100 mechanism gate;
  \item the frozen positive-control construction was infeasible for the chosen
  model/decoding contract; and
  \item the preregistered causal and confirmatory questions remain unanswered.
\end{enumerate}

The campaign does \emph{not} establish a general equivalence of TRL and the
reference, a harmful training effect from native semantics, or an advantage from
a controller.

\section{Discussion}

\paragraph{Optimizer efficiency is not sample efficiency.}
For fixed $G$, Eq.~\eqref{eq:boundary-cost} separates group count from rollout
count; Eq.~\eqref{eq:all-g-cost} gives an all-$G$ lower bound for complete
groups.  Neither result analyzes an adaptive stopping policy.  A publishable
controller result should report both contrast yield and charged rollout cost.

\paragraph{Conformance is not sameness.}
Native frameworks may intentionally choose different semantics.  The scientific
requirement is to bind a claim to the semantics that actually ran and to avoid
attributing a difference to an algorithm when row selection or reduction
changed.  ``Material difference'' is a routing signal for investigation, not a
defect label.

\paragraph{Negative and incomplete outcomes are informative.}
The filtered regime's infeasibility falsifies a property of the experimental
construction, not the causal hypothesis itself.  Silently lowering the CV gate
would create a different study.  Preserving the failure prevents an attractive
but unsupported positive-control narrative.

\paragraph{Operational consequence.}
An outcome-only controller should expose uncertainty rather than convert a
homogeneous group directly into ``retry'' or ``drop.''  The next action can be a
small information-gathering probe: reverify an answer, inspect an execution
receipt, or sample a bounded number of additional trajectories.  The probe is
justified only when its expected decision value exceeds its cost.

\section{Limitations}

The empirical study uses one 1.7B model family, one completed task regime, two
partially represented seeds, and a short 100-step horizon.  The result is not a
benchmark of framework quality.  The S1 cases are deliberately small and do not
prove end-to-end equivalence.  The action-reversal theorem abstracts away model
updates between retries, correlated rollouts, continuous rewards, and changing
verification policies.  Its purpose is to prove insufficiency under transparent
assumptions, not to prescribe a universal controller.

Finally, the current repository contains mutable working-tree development
around immutable acceptance artifacts.  Reproduction should use the source
hashes and artifact commits in receipts, not assume that a directory name alone
identifies a run.

\section{Conclusion}

A homogeneous reward group certifies missing within-group contrast; it does not
identify why the contrast is missing or which intervention is optimal.  Near a
binary reward boundary, larger fixed groups reduce the expected number of
groups to contrast, while complete-group sampling obeys a $1/\epsilon$ rollout
lower bound.  In the restricted stop-versus-batch problem, two compatible
latent states demand opposite actions and yield positive two-point minimax
regret.

These facts make semantic and evidential discipline part of the algorithm.  A
fail-closed conformance chain can show what objective, gradient, and update
actually ran; immutable receipts can distinguish accepted evidence from failed,
descoped, pending, and unrun work.  In our campaign, that discipline exposes a
failed balanced mechanism gate and an infeasible positive-control construction,
not the preregistered causal headline.

\section*{Reproducibility Statement}

The content-addressed \path{review_bundle.zip} contains the paper source/PDF,
S1 source and receipts, the frozen preregistration, all r4-2 acceptance records,
campaign state, recovery evidence, and frozen source-provenance archives.  Its
internal \path{MANIFEST.sha256} covers every payload file; the outer digest is
in \path{REVIEW_BUNDLE.sha256}.  After extraction, run
\begin{verbatim}
python3 verify_claims.py --repo-root repository
\end{verbatim}
to run formula sanity checks and verify hashes plus internal receipt invariants,
including all 600 stored gradient relations and six evaluation/compute ledgers.
The script performs no network access and does not regenerate gradients,
predictions, or training data from private checkpoints.  Exact coverage and
limitations are documented in \path{CLAIM_AUDIT.md} and
\path{REVIEW_BUNDLE.md}.

\bibliographystyle{plainnat}
\bibliography{flagship}

\appendix

\section{Exact Group-Size Definition and Bounds}

For target homogeneous probability $\tau\in(0,1)$, define the exact integer
choice for $p\in(0,1)$
\begin{equation}
 G^\star(p,\tau)=\min\{G\geq2:f_G(p)\leq\tau\}.
\end{equation}
At $p\in\{0,1\}$ define $G^\star(p,\tau)=+\infty$.  For $p\in(0,1)$, let
$q=\max(p,1-p)\in[1/2,1)$.  Since
$q^G\leq f_G(p)\leq2q^G$, the exact choice obeys the integer bracket
\begin{equation}
 \max\!\left\{2,\left\lceil\frac{\log\tau}{\log q}\right\rceil\right\}
 \leq G^\star(p,\tau) \leq
 \max\!\left\{2,\left\lceil\frac{\log(\tau/2)}{\log q}\right\rceil\right\}.
\end{equation}
We recommend exact integer search rather than dropping the smaller term in
Eq.~\eqref{eq:degeneracy}.

\section{Decision-Reversal Numerical Witness}

The verification script instantiates Proposition 3 with $q=0.1$, $m=8$,
$c=0.02$, and $B=1$.  Retry has incremental utility $-0.16$ in $E_0$ and
$0.40953279$ in $E_q$.  The resulting minimax-regret lower bound is
$0.115050876$.  These numbers are a sanity-check witness, not an empirical
estimate for the campaign.

\section{Frozen S1 Case Matrix}
\label{app:cases}

Table~\ref{tab:s1-cases} enumerates the 14 intended cases executed separately
against both stack adapters.  ``Selected'' means rows $0,1,3$ are retained and
all other completion masks are zero.  The fixture definitions and exact receipt tensors
are in the review bundle.

\begin{table}[h]
  \centering
  \caption{Complete intended S1 case inventory.}
  \label{tab:s1-cases}
  \small
  \begin{tabularx}{\linewidth}{lXl}
    \toprule
    Arm & Fixtures & Targeted semantic \\
    \midrule
    GRPO & base, all-wrong, all-correct, graded, translated, low-clip
      & standardization, zero branch, translation, lower clip \\
    DAPO & DAPO-clip, low-clip & asymmetric upper/lower clipping \\
    GSPO & base, zero-mask & sequence ratio and masking \\
    DrGRPO & base, translated & unscaled centered reward and reduction \\
    AERO & posterior & posterior advantage on homogeneous rewards \\
    GRPO selected & base; rows 0,1,3 & active-row selection \\
    \bottomrule
  \end{tabularx}
\end{table}

Native TRL is materially different from the reference on GRPO, DAPO, GSPO,
and DrGRPO and does not expose the tested AERO objective.  Native verl is
materially different on GRPO; DAPO, GSPO, DrGRPO, and AERO are not exposed by
the pinned kernel.  ``Material'' means the stored native trace differs from
this reference under the frozen tolerances, not that the framework is defective.
The shared controller cases are the Cartesian product
\{static-8, static-16, symmetric-ZVF, failure-only, boundary-aware,
full-triage\} $\times$ \{all-wrong, all-correct, mixed, noisy, missing,
delayed\}.

\section{Complete r4-2 Job Crosswalk}

Table~\ref{tab:crosswalk} partitions all 31 manifest IDs.  The verifier expands
and checks the exact identifiers, so no unlisted filename glob can change the
reported decomposition.

\begin{table}[h]
  \centering
  \caption{Exhaustive planned-job crosswalk.}
  \label{tab:crosswalk}
  \small
  \begin{tabularx}{\linewidth}{Xrl}
    \toprule
    Manifest ID set & Count & Disposition \\
    \midrule
    \texttt{preflight\_\_a100\_stack\_smoke} & 1 & accepted \\
    balanced corpus, seeds 11/23/37 & 3 & accepted \\
    balanced scientific, seed 11, epsilon/reduction & 2 & accepted \\
    balanced scientific, seed 11, intended/native & 2 & failed infrastructure \\
    balanced scientific, seed 23, four conditions & 4 & accepted \\
    balanced scientific, seed 37, four conditions & 4 & pending quota \\
    filtered corpus, seed 11 & 1 & failed validation \\
    filtered corpus, seeds 23/37 & 2 & contract-infeasible \\
    filtered scientific, three seeds, four conditions & 12 & contract-infeasible \\
    \midrule
    Total & 31 & 10/14/2/1/4 state counts \\
    \bottomrule
  \end{tabularx}
\end{table}

\section{Evidence Status Matrix}

\begin{table}[h]
  \centering
  \caption{Boundary between current evidence and future work.}
  \begin{tabularx}{\linewidth}{lXl}
    \toprule
    Item & Evidence & Status \\
    \midrule
    Binary homogeneous probability & Analytic identity; prior work & Sourced \\
    First-contrast and boundary cost & Proofs plus executable checks & Proved \\
    Two-action batch-retry reversal & Proof plus numerical witness & Proved \\
    S1 intended conformance & Frozen source/fixture receipt & Verified \\
    Balanced gradient relations & Six final acceptance records & Verified \\
    Filtered positive control & Frozen CV gate cannot be met & Infeasible \\
    Causal training effect & Needs both regimes and all seeds & Not established \\
    Controller benefit & No matched intervention study & Not run \\
    \bottomrule
  \end{tabularx}
\end{table}

\end{document}
